% ==============================================================================
% sections/introduction.tex
% Created: 2026-05-08
% ==============================================================================

Distinguishing chronic from transitory unemployment is central to diagnosing the
nature of joblessness and designing effective labour market policy. The answer
determines which policy lever to pull. In economies with high labour market
churn---where workers frequently move into and out of employment---the policy
challenge is to protect workers against frequent spells of unemployment, and
policies favouring stronger worker protections and enforcement are appropriate
\citep{aghion2006}. By contrast, if employment is highly persistent, unemployment
reflects a failure of job creation rather than a failure of retention; the
policy prescription shifts toward lowering hiring costs and removing barriers to
mobility \citep{nattrass2000}. Because transition matrices are the primary tool
for assessing the degree of labour market churn, credible measurement of
entry and exit rates is not merely an academic exercise---it is a prerequisite
for coherent policy design.

Labour force surveys are the workhorse for measuring employment transitions.
Yet survey errors are ubiquitous, and classification error in employment status
is especially damaging for transition estimates. The effect of misclassification
on flows is multiplicative: a single transient error at wave $t$ makes a
chronically non-employed worker appear to have found a job and then lost it,
fabricating two transitions that did not occur. Even low misclassification rates
therefore produce severe upward bias in estimated flows
\citep{porteba1986,singh1995,meyer1988}. The same worker also appears to
exhibit negative duration-dependence---a recent ``entrant'' who immediately
``exits''---when in reality the worker never moved. Classification error thus
simultaneously overstates churn and attenuates persistence.

This paper develops a structural estimator that jointly identifies true
employment transition rates and misclassification probabilities using only three
consecutive waves of panel data and no reinterview or administrative validation.
The estimator is based on an Expectation--Maximisation (EM) algorithm that
exploits a simple and powerful identification insight: if the latent true
employment process follows a first-order Markov chain but observed employment is
measured with error, then the observed three-wave distribution will systematically
violate the Markov property in a way that separately pins down the transition
rates and the misclassification probability. The EM algorithm converts this
insight into a sequence of closed-form updates, making estimation fast,
numerically reliable, and easily extended.

A large literature has documented the importance of classification error in
employment surveys. The pioneering work of \citet{porteba1986}, using reinterview
data for the United States Current Population Survey (CPS), found that correcting
for classification error reduced the monthly probability of job loss from 5\%
to 1.9\% and job entry from 7.2\% to 2.6\%. \citet{chua1987} found, using the
same data, a 6\% chance of misclassifying a truly unemployed individual as
employed. \citet{singh1995} estimated misclassification probabilities of less
than 1\% in Canada, yet showed that even such small error rates produce
substantial upward bias in transition estimates. \citet{abowd1985} introduced the
error-correction matrix approach that underlies much of this work. Reinterviews,
however, are costly, rare in the lower income countries where survey data are 
more prone to error \citep{ashenfelter1986}. Moreover, these auxiliary data are
themselves imperfect---respondents may change status between the original and
re-interview survey, or modality differences between telephone and in-person
interviews may introduce systematic bias (CITE LSMS NIGERIA PAPER?).

Motivated by these limitations, several estimators have been proposed that
identify misclassification from panel data alone.
\citet{biemer2000} introduced Markov latent class analysis (MLCA) to the
three-wave CPS data and showed that the resulting estimates closely match
reinterview-based corrections. Their approach imposes a first-order Markov
assumption on the true process and a local independence assumption on the
errors. \citet{feng2013} use an eigenvalue decomposition of the response
classification matrix to estimate misclassification rates and correct unemployment
stock estimates for the United States, finding that official unemployment is
understated by approximately 2.1 percentage points. \citet{shibata2019} develops
a hidden Markov model (HMM) formulation that accommodates individual-level
heterogeneity in the latent process, returning to maximum likelihood estimation
rather than the eigenvalue approach. Our estimator builds on this tradition but
contributes to it in four ways.

First, we derive and implement a modular EM algorithm that yields closed-form
M-step updates for the baseline model, making estimation scalable to large
samples while allowing straightforward extension. Second, we develop a
generalized EM (GEM) extension that allows observable worker heterogeneity via
probit transition probabilities and a finite mixture model (FMM) for unobservable
heterogeneity. Both of these can mimic Markov violations---a worker who always
has high exit probability looks like a Markov violation if heterogeneity is
ignored---and the GEM extension ensures these alternatives cannot masquerade as
misclassification. Third, we extend the model to a second-order Markov
specification using four waves of data, directly testing and controlling for
duration dependence in the latent process. Fourth, we develop a novel
tenure-contamination model in which reported job tenure provides a structural
identification channel: if employment status is misclassified at an intermediate
wave, the tenure reported at flanking waves must be inconsistent with the conditionally 
deterministic tenure accumulation. The model separates tenure contamination (random recall
error in tenure reporting) from state misclassification, providing an
independent corroboration of the baseline result.

We apply the estimator to the South African Quarterly Labour Force Survey (QLFS),
a rotating panel survey collected quarterly since 2008. South Africa provides
an ideal empirical setting for two reasons. First, it has persistently high
unemployment rates---around 30\%---alongside a vigorous academic debate about
the nature of joblessness, with important policy implications for the design of
labour market interventions \citep{banerjee2008,nattrass2000,burgervonfintel2014}.
Second, there is well-documented evidence of measurement problems in the QLFS
\citep{burger2016,lechtenfeld2014}, and reinterview data are unavailable,
making a structural approach attractive.

Using the QLFS data, we document a striking empirical puzzle: three-month
employment entry and exit rates, estimated from raw transition matrices, are
approximately 8\%, implying six-month rates of roughly 15\% under the Markov
assumption. But the observed six-month transition rates---computed directly from
four-wave panels---are only about 10\%, roughly one-third lower than the implied
value. This divergence is inconsistent with a first-order Markov process on
the observed employment status, but is exactly what one would expect if the
latent true process satisfies the Markov assumption and observed employment is
measured with transient error.

Our baseline EM estimator recovers a misclassification probability of
approximately 3\%, which implies true quarterly entry and exit rates of
approximately 3\%---a threefold correction relative to the naive transition
matrix. The symmetric and asymmetric misclassification specifications produce
virtually identical results, as do the stationarity-restricted and
free-initial-employment-probability versions. The estimated misclassification
probability remains stable---between 2\% and 3\%---across extensions that 
controll for observed heterogeneity (age, education, gender, race, contract
type), unobserved heterogeneity (two-type finite mixture), and second-order
Markov dynamics. When we allow the misclassification probability to vary with
survey response reliability---proxied by inconsistencies in reported age
progressions or education levels---we find that the baseline misclassification
rate is approximately 2.2\% for reliable respondents but rises to over 17\%
for those whose age changes by an implausible amount between waves. This
mechanism evidence reinforces the interpretation that the estimated
misclassification reflects genuine survey error, not a modelling artefact.
ADD HERE ABOUT THE TENURE CONTAMINATION MODEL. Stricter wave-to-wave matching criteria, which reduce the likelihood of
including mismatched observations, consistently lower both the measured
transition rates and the estimated misclassification probability, providing
a further design-based robustness check.

The findings have substantive implications for the characterisation of South
African unemployment. If true quarterly entry and exit rates are roughly 3\%,
the average employment spell is on the order of several years, not several
quarters. Employment in South Africa is considerably more persistent than
previously characterised by studies using raw transition matrices
\citep{banerjee2008,cichello2014,kerr2018}. This is more consistent with a
labour market in which unemployment reflects scarce job creation rather than
high churn, pointing toward supply-side and job-creation policies rather than
policies targeted at worker protection or flexible hiring.

The remainder of the paper is organised as follows. Section~\ref{sec:data}
describes the QLFS data and documents the empirical puzzle motivating the
approach. Section~\ref{sec:methodology} presents the EM estimator and its
extensions. Section~\ref{sec:results} reports the main results.
Section~\ref{sec:alternatives} reports the alternative specifications.
Section~\ref{sec:conclusion} concludes. Full technical derivations for the
baseline, AR(2), and tenure contamination models are provided in
Appendices~\ref{app:baseline}, \ref{app:ar2}, and~\ref{app:tenure}.
